\section{Introduction}
\label{sec:introduction}

Vision-language systems and embodied agents are increasingly tested in
simulated 3D scenes, where they must bind language to objects, views, and
actions. These evaluations treat the rendered scene as fixed input, yet the
view is also the measurement interface. When scenes cross shader systems for
large runs, sharing, viewer export, or robotics reuse, that step can become
a hidden task factor.

In Isaac Sim, many photorealistic assets use NVIDIA's MDL shaders (Material
\mbox{Definition} Language). Portable USD and glTF pipelines usually expect
UsdPreviewSurface or PBR materials. That step helps reuse, but it can also
shift color, gloss, lighting, or texture cues. For grounding benchmarks, visual
match is only the first question. The next question is whether a model still
answers the same object query and points to the same object after the material
changes.

We therefore ask which checks must hold as converted assets move into grounding
and stack checks, because a single proxy score is not enough. In
GRScenes~\cite{Wang2024GRUtopia}, runs combine large USD scenes, object goals,
VLM prompts, and VLN routes such as \mbox{InternNav} and
\mbox{DualVLN}~\cite{Wei2026Ground}.

\begin{figure*}[t]
\centering
\includegraphics[width=\textwidth]{figures/fig_acl_method_chain_imagegen_v18.png}
\caption{Evidence-chain overview. We start from
USD scenes with MDL materials, convert them to composition-preserving noMDL
counterparts, compare matched renders for VLM grounding, then check selected
official InternNav rollouts and paired metrics. This
panel is a schematic roadmap; later render, material-effect, and navigation
figures carry the reported measurements.}
\label{fig:acl_method_chain}
\end{figure*}

Figure~\ref{fig:acl_method_chain} gives the schematic roadmap. The protocol has
four checks. First, four Isaac Sim assets test visual and feature match on noMDL
outputs. Next,
GRScenes grounding keeps source scenes unchanged, uses scratch noMDL rewrites,
matched renders, target boxes, blind visual QA, and local Gemma4 and
\mbox{Qwen2.5-VL} probes. The clean pool has 15 pairs; the frozen stress pool
has 30 pairs and tests objects across color, gloss, or light
shifts. Third, one controlled comparison covers source MDL, noMDL, and NVIDIA
outputs in checked bins. Fourth, an official \mbox{InternNav} 99-episode run
tests noMDL stack entry. Together, these checks make the evaluation scope
explicit.

We make four contributions that turn a hidden preprocessing step into an
auditable grounding measurement protocol:
\begin{enumerate}
    \item We treat material conversion as a controlled VLM-grounding
    intervention rather than routine asset cleanup.
    \item We specify a controlled perturbation protocol with immutable
    inputs, scratch conversion, matched renders, projection QA, prompt
    formats, and coordinate scoring.
    \item We instantiate the protocol on GRScenes with a clean visual-QA
    pilot and a frozen 30-pair target-centered stress set, showing backend
    differences in answer, point, and coordinate checks.
    \item We keep material and stack checks limited to an NVIDIA material audit,
    99 KuJiaLe VLN episodes, and repeated load/render checks. These checks
    report entry and stability only for the tested settings; broader benchmark and
    speed comparisons remain outside this study.
\end{enumerate}
